% nif_artgrav: artificial gravity in a rotating habitat (Ch 10, fig:artgrav)
% Left: the mechanism, in the ground (inertial) frame.  The only real force on the
% astronaut is the floor's push N = m omega^2 r, directed INWARD (towards the hub).
% Near the hub r ~ 0 and there is no floor to stand on: the crew float.
% Right: the radius needed for g_eff = omega^2 r = g = 9.8 m/s^2 at spin rate n (rpm),
%   r = g / (2 pi n / 60)^2:  n = 2 -> 223 m,  n = 3 -> 99 m,  n = 5.5 -> 29.5 m.
% The Endurance (Interstellar) ring, r = 32 m at 5.5 rpm, gives 1.08 g, so its point sits
% just above the 1 g curve.  Comfort bands as usually quoted: below 2 rpm comfortable,
% 2 to 6 rpm tolerable with adaptation, above 6 rpm too fast.  The curve starts at
% 0.55 rpm, where r = 2955 m.
\begin{tikzpicture}[x=1cm,y=1cm]
  % ================================================================= habitat
  \begin{scope}
    \def\Ro{2.25} \def\Ri{1.85} \def\Rh{0.36}
    \draw[outline,fill=cMoon,even odd rule] (0,0) circle[radius=\Ro] circle[radius=\Ri];
    \foreach \a in {22.5,67.5,...,337.5}
      \draw[cGrey,thin line] (\a:\Rh) -- (\a:\Ri);
    % radius from the axis to the floor (starts at the centre, under the hub)
    \draw[thin line,cSteel,-smallhead] (0,0) -- node[pos=0.6,anchor=south west,cSteel,inner sep=1pt] {$r$} (-30:\Ri);
    \draw[outline,fill=cMoon] (0,0) circle[radius=\Rh];
    \node[font=\footnotesize] at (0,0) {hub};
    % spin
    \draw[rotation] (60:\Ro+0.3) arc[start angle=60,end angle=120,radius=\Ro+0.3];
    \node[anchor=south] at (90:\Ro+0.35) {$\omega$};
    % astronaut standing on the floor of the rim, head towards the hub
    \begin{scope}[shift={(0,-\Ri)}]
      \draw[body,rounded corners=1.5pt] (-0.09,0.02) rectangle (0.09,0.52);
      \draw[body] (0,0.64) circle[radius=0.11];
    \end{scope}
    \draw[force] (-0.22,-\Ri) -- ++(0,0.9) node[anchor=south,inner sep=2pt] {$N$};
    % a crew member floating near the hub
    \begin{scope}[shift={(-0.95,0.05)},rotate=-60]
      \draw[body,rounded corners=1.5pt] (-0.07,-0.2) rectangle (0.07,0.2);
      \draw[body] (0,0.3) circle[radius=0.085];
    \end{scope}
    \draw[leader] (-1.15,0.2) -- (-2.5,1.25);
    \node[anchor=east,font=\footnotesize,align=right,inner sep=1pt] at (-2.5,1.25) {floats\\ near the\\ hub};
    \draw[leader] (-0.3,-1.35) -- (-2.5,-1.35);
    \node[anchor=east,font=\footnotesize,align=right,inner sep=1pt] at (-2.5,-1.35) {floor\\ pushes\\ inward};
  \end{scope}
  % ================================================================= design plot
  \begin{axis}[bookplot,
      at={(4.1cm,-2.0cm)}, anchor=origin,
      width=8.2cm, height=5.0cm, scale only axis=false,
      ymode=log, log basis y=10,
      xmin=0, xmax=7.5, ymin=8, ymax=3000,
      xtick={0,1,...,7}, ytick={10,100,1000},
      yticklabels={$10$,$100$,$1000$},
      xlabel={spin rate (rpm)},
      ylabel={radius for $1\,g$ (m)},
      domain=0.55:7.5]
    % comfort bands
    \fill[cFor!10]   (axis cs:0,8) rectangle (axis cs:2,3000);
    \fill[cAmber!22] (axis cs:2,8) rectangle (axis cs:6,3000);
    \fill[cCrim!10]  (axis cs:6,8) rectangle (axis cs:7.5,3000);
    \node[cFor,font=\footnotesize,anchor=south] at (axis cs:1,9.5) {comfortable};
    \node[font=\footnotesize,anchor=south] at (axis cs:4,9.5) {tolerable with adaptation};
    \node[cCrim,font=\footnotesize,anchor=north] at (axis cs:6.75,2700) {too fast};
    % r = g / omega^2
    \addplot[cPrim] {9.8/(2*pi*x/60)^2};
    \foreach \n in {2,3}{
      \pgfmathsetmacro\rr{9.8/(2*pi*\n/60)^2}
      \edef\temp{\noexpand\node[point mass,fill=cPrim] at (axis cs:\n,\rr) {};}\temp}
    \node[cPrim,anchor=south west,inner sep=2pt] at (axis cs:2,223.4) {$223$ m};
    \node[cPrim,anchor=south west,inner sep=2pt] at (axis cs:3,99.3) {$99$ m};
    % the Endurance
    \node[dot,cCrim] at (axis cs:5.5,32) {};
    \draw[leader,cCrim] (axis cs:5.45,36) -- (axis cs:4.9,230);
    \node[cCrim,anchor=south,align=center,inner sep=2pt] at (axis cs:4.9,230)
      {\emph{Endurance}:\\ $32$ m, $5.5$ rpm};
  \end{axis}
\end{tikzpicture}
